\documentclass[leqno, a4paper, 12pt]{jsarticle}
\usepackage{amssymb}
\usepackage{amsthm}
\usepackage{amsmath}
\usepackage{comment}
	%可換図式用
		%\usepackage[all]{xy}
	%重くなるので使わいときはコメント化しておく。
%定理環境 thmはあまり使わずpropをなるべく使う。
%主張は基本的に証明の中で使い、そのため番号を付けない。
%注意環境を付けてみた。
\newtheorem{thm}{定理}
\newtheorem{prop}[thm]{命題}
\newtheorem{cor}[thm]{系}
\newtheorem{lem}[thm]{補題}
\newtheorem{df}[thm]{定義}
\newtheorem{exam}[thm]{例}
\newtheorem{fact}[thm]{事実}
\newtheorem*{claim}{主張}
\newtheorem*{rmk}{注意}
\renewcommand{\proofname}{{\bf 証明}}
%矢印たち。
%\toは\rightarrowと同じ。\getsは\leftarrowと同じ。同値は\iff
%上向き矢はuparrow逆はdownarrow
\newcommand{\To}{\Longrightarrow}
\newcommand{\Gets}{\Longleftarrow}
%黒板太字
\newcommand{\nn}{\mathbb{N}}
\newcommand{\zz}{\mathbb{Z}}
\newcommand{\qq}{\mathbb{Q}}
\newcommand{\rr}{\mathbb{R}}
\newcommand{\cc}{\mathbb{C}}
%無限大
\newcommand{\mgn}{\infty}
%ギリシャン
\newcommand{\dpst}{\displaystyle}
\newcommand{\ep}{\varepsilon}
\newcommand{\al}{\alpha}
\newcommand{\be}{\beta}
\newcommand{\de}{\delta}
\newcommand{\la}{\lambda}
\newcommand{\La}{\Lambda}
\newcommand{\ga}{\gamma}
\newcommand{\lil}{\la\in \La}
%商集合
\newcommand{\qt}[2]{#1/\! #2}
%enumerate環境
\renewcommand{\labelenumi}{{\rm (\arabic{enumi})}}
%以降alignじゃなくてenumerate を使え。
%なんか演算
\DeclareMathOperator{\card}{card}
\DeclareMathOperator{\supp}{supp}
%なんか演算２
\DeclareMathOperator{\bd}{Bdry}
\DeclareMathOperator{\cl}{CL}
\DeclareMathOperator{\nai}{Int}
\DeclareMathOperator{\di}{diam}

%差集合は\setminusで統一する。等号の否定は\neq
%真部分集合は\subsetneqq　偏微分は\partial
%ディスプレイスタイル。\dfracでディスプレイ分数
%空集合は\Oを使う！！！！！！！！！！！！

\title{Lebesgue--Stieltjes Measure}
\author{ナンブキトラ}
\date{\today}
			\begin{document}
\maketitle



この文章ではLebesgue--Stieltjes測度を構成してゆく．
つまり以下の定理を証明する．
\begin{thm}
$\phi:\rr\to \rr$を右連続な非減少関数とする．このとき，
$\rr$上の測度$\mu_{\phi}$で，
任意の$a, b\in [-\infty, \infty]$で$a<b$を満たすものに対して
\[
\mu_{\phi}((a, b])=\phi(b)-\phi(a)
\]
となるものが存在する．
\end{thm}


\section{Riesz--Markov--Kakutaniによる方法}

以下の内容については\cite{RMK1}や
\cite{RMK2}を読んでも良い．

\begin{df}
$X$を局所コンパクトハウスドルフ空間とする．
このとき$X$上の測度$\mu$がRMK測度であるとは以下の条件を満たすときにいう
\begin{enumerate}
\item $\mu$は$[0, \infty]$に値をとるボレル測度である．
\item $X$の任意のコンパクト集合$K$について$\mu(K)<\infty$
\item $X$の任意のボレル集合は外部正則である．つまり
$\mu(E)=\inf\{\mu(V)\ |\ V\in \mathfrak{O}_X\ E\subset V\}$が成り立つ．ここで$\mathfrak{O}_X$は$X$の開集合系である．
\item $E$が開集合もしくは$\mu(E)<\infty$ならば$E$は
コンパクト正則である．つまり，
$\mu(E)=\sup\{\mu(K)\ |\ K\in Cpt(X) \ K\subset E\}$
が成り立つ．ここで$Cpt(X)$は$X$のコンパクト集合全体である．
\end{enumerate}

$X$上のRMK測度全体を$M(X)$と書くことにする
$X$上の符号付き測度で二つのRMK測度の差で表されるもの全体を$S(X)$と書くことにする．
\end{df}


\begin{df}
$X$を位相空間とする．
このとき
$C(X)$で$X$上の実連続関数全体の集合を表す・
$C_{b}(X)$で$X$上の有界実連続関数全体の集合を表す．
また，$X$が距離空間で距離関数$d$を備えているとき，
$C_{u}(X)$で$X$上の($d$に関して)一様連続な実関数全体を表す．$d$を明示的に書いていないが，以下では紛れがないので心配は無用である．
また，以下では$C_{b}(X)$には$\sup$ノルムが備わっているとする．
\end{df}

以下の定理はRiesz--Markov--Kakutaniの表現定理の一つの形態である．
\begin{thm}\label{thm:rmkdual}
$X$をコンパクトハウスドルフ空間とする．
このとき
\[
C_b(X)^*=C(X)^*=S(X). 
\]
である．
より詳しく述べると，$\mu\in C_{b}(X)^{*}$に対して
二つのRMK測度$\alpha$と$\beta$が（差を除いて）一意的に対応して，
任意の$f\in C_{b}(X)$に対して
\[
\mu(f)=\int_{X}f \ d\alpha -\int_{X}f\ d\beta
\]
が成り立つ．
特に非負な$\mu\in C_{b}(X)^{*}$つまり
$f\ge 0$ならば$\mu(f)\ge 0$を満たすような$\mu$
は通常の意味での測度（非負の値しか取らない測度）
に対応する．
\end{thm}



\begin{df}
$\phi:\rr\to \rr$を右連続な非減少関数とする.
このとき，
任意の$n\in \nn$に対して
$L_{\phi, n}:C_{c}(\rr)\rr$
を以下のように定義する．
\[
L_{\phi, n}(f)=\sum_{i=-\infty}^{\infty}f\left(\frac{i}{2^n}\right)\left(\phi\left(\frac{i}{2^n}\right)-\phi\left(\frac{i-1}{2^n}\right)\right)
\]
と定義する．$f$の台がコンパクトなので，これは有限和である．
\end{df}

\begin{lem}
$\phi:\rr\to \rr$を右連続な非減少関数とする.
任意の$f\in C_c(\rr)$に対して
$\lim_{n\to \infty}L_{\phi, n}(f)$が存在する．
\end{lem}
\begin{proof}
$M>0$を$\supp (f)\subset [-M, M]$を満たすものとし，
$C=\phi(2M)-\phi(-2M)$とする．
任意に$\ep>0$を与える．
このとき$f$は一様連続なので，ある$\delta$
が存在して，
$|x-y|<\delta$ならば$|f(x)-f(y)|<\ep/C$を満たす．
$N\in \nn$を$2^{-N}<\delta$とすると，
$N<m<n$をみたす$n, m\in \nn$について，
$\{a_i\}\subset \rr$
を
$0\le k<2^{n-m}$ならば
$a_{i+k}=a_i$
とし，$a_{k(2^{n-m})}=k$
と定義すると，
\[
L_{\phi, m}(f)
\sum_{i=-\infty}^{\infty}f\left(\frac{a_i}{2^n}\right)
\left(\phi\left(\frac{i}{2^n}\right)-\phi\left(\frac{i-1}{2^n}\right)\right)
\]
となるから，
\begin{align*}
|L_{\phi, n}(f)-L_{\phi, m}(f)|
&=
\sum_{i=-\infty}^{\infty}
\left|f\left(\frac{i}{2^n}\right)-f\left(\frac{a_i}{2^{-n}}\right)\right|
\left(\phi\left(\frac{i}{2^n}\right)-\phi\left(\frac{i-1}{2^n}\right)\right)\\
&<
\frac{\ep}{C}\sum_{i=-\infty}^{\infty}\left(\phi\left(\frac{i}{2^n}\right)-\phi\left(\frac{i-1}{2^n}\right)\right)
\le 
\ep
\end{align*}
を得る．
つまり，$\{L_{\phi, n}(f)\}_{n\in \nn}$
はコーシー列なので，
$\lim_{n\to \infty}L_{\phi, n}(f)$が存在する．
\end{proof}


\begin{df}
$\phi:\rr\to \rr$を右連続な非減少関数とする.
このとき，
$L{\phi}: C_c(\rr)\to \rr$を
\[
L_{\phi}(f)=\lim_{n\to \infty}L_{\phi, n}(f)
\]
と定義する．
\end{df}


\begin{lem}\label{lem:senkei}
$\phi:\rr\to \rr$を右連続な非減少関数とする.
このとき$L_{\phi}$
は正値線形形式である．
\end{lem}
\begin{proof}
各$L_{\phi, n}$が正値線形形式だからその極限で定義される
$L_{\phi}$もそうである.

\end{proof}

\begin{df}
$\phi:\rr\to \rr$を右連続な非減少関数とする.
上の補題 \ref{lem:senkei} とRMKから，
$L_{\phi}$に対応する
RMK測度が存在する．
これを$\mu_{\phi}$と書くことにする．
もちろん
\[
L_{\phi}(f)=\int_{\rr}f \ d\mu_{\phi}
\]
が成り立つ．
\end{df}


\begin{prop}
$\phi:\rr\to \rr$を右連続な非減少関数とする.
このとき
任意の$a, b\in [-\infty, \infty]$で$a<b$を満たすものに対して
\[
\mu_{\phi}((a, b])=\phi(b)-\phi(a)
\]
が成り立つ．
\end{prop}
\begin{proof}
まず最初に, $a, b$が
$s, t\in \zz$で$N\in \zz_{\ge 0}$
を使って
$a=s/2^N$, 
$b=t/2^N$
と表される場合に等式を証明する．
$n\in \zz_{\ge 0}$を$N<n$として，
$c_n=(s2^{n-N}+1)/2^N$, 
$d_n=(t2^{n-N}+1)/2^N$
とおく．
このとき$f_n\in C_c(\rr)$を以下のように定義する．
\begin{align*}
f_n(x)=
\begin{cases}
\frac{1}{c_n-a}(x-a) & \text{if $x\in [a, c_n]$}\\
1 & \text{if $x\in [c_n, b]$}\\
\frac{-1}{d_n-b}(x-d_n) & \text{if $x\in [b, d_n]$}\\
0 & \text{o. w. }
\end{cases}
\end{align*}
すると，$n\to \infty$のとき$f_n\to \chi_{(a, b]}$
なので，
\[
\mu_{\phi}((a, b])=\lim_{n\to \infty}\int_{\rr}f_n\ d\mu_{\phi}
=
\lim_{n\to \infty}L_{\phi}(f_n)
\]
である．
$m>n$として，
$L_{\phi, m}(f_n)$を計算する．
\[
|L_{\phi, m}(f_n)-(\phi(b)-\phi(a))|\le (\phi(d_n)-\phi(b))+(\phi(c_n)-\phi(a))
\]
となるので特に
\[
|L_{\phi}(f_n)-(\phi(b)-\phi(a))|\le (\phi(d_n)-\phi(b))+(\phi(c_n)-\phi(a))
\]
である．
$\phi$は右連続なので，$n\to \infty$のとき右辺は$0$に近づくから，
\[
\lim_{n\to \infty}L_{\phi}(f_n)=\phi(b)-\phi(a)
\]
よって
\[
\mu_{\phi}((a, b])=\phi(b)-\phi(a)
\]
がわかる．$a, b$が一般の場合には，測度論を使えばわかる．
つまり，
$a_{n}$と$b_{n}$を$s/2^{N}$ $(s\in \zz, N\in \zz_{\ge 0})$の形で表される数で
$a<a_{n}$と$b<b_{n}$を満たし
$\lim_{n\to \infty}a_{n}=a$と
$\lim_{n\to \infty}b_{n}=b$
を満たすものとする．
すると
$\bigcup_{n\in \zz_{\ge 0}}\bigcap_{m\in \zz_{\ge 0}}(a_{n}, b_{m}]=(a, b]$
なので，有限測度に関する極限公式と$\phi$の右連続性を使えば
$\mu_{\phi}((a, b])=\phi(b)-\phi(a)$
がわかる．
\end{proof}

\section{分位関数を使う方法}
この節では，
Lebesgue--Stieltjes測度の存在を
Quantile function（分位関数）を使って証明する．
ただし，通常のLebesgue測度の存在は前提とする点には注意しよう．

\begin{df}
$\phi$を$\rr$上の右連続な単調増加関数とする．
このとき，
$I_{\phi}=[\inf\phi, \sup\phi]\cap \rr$と定義する
（ここでわざわざ$\rr$との共通部分をとっているのは，
$\sup$や$\inf$が$\infty$や$-\infty$の値を取る場合を見越してのことであり，そうでない場合には余計なものである）．
$\phi$の\textbf{分位関数}$Q_{\phi}: I_{\phi}\to \rr$を以下のように定義する．
\[
Q_{\phi}(w)=\inf\{\, x\in \rr\mid w\le \phi(x)\, \}
\]
分位関数のことを\textbf{一般化した逆関数}と呼んだりもする．
\end{df}



\begin{lem}
任意の$w\in I_{\phi}$と$x\in \rr$について
$w\le \phi(x)$と$Q_{\phi}(w)\le x$
は同値である．
\end{lem}
\begin{proof}
まず$w\le \phi(x)$を仮定する．
このとき$Q_{\phi}$の定義から
$Q_{\phi}(w)\le x$である．

逆に$Q_{\phi}(w)\le x$を仮定する．
このとき$\phi$の右連続性から
$\{\, z\in \rr\mid w\le \phi(z)\, \}$
は閉集合であるから
この集合の$\inf$である$Q_{\phi}(w)$も
この集合に属する．よって
$w\le \phi(Q_{\phi}(w))$である．
今，$Q_{\phi}(w)\le x$と仮定しているので，
$\phi$の単調性から$w\le \phi(x)$である．
\end{proof}

\begin{cor}
$x\in \rr$を固定すると
\[
\{\, w\in I_{\phi}\mid Q_{\phi}(w)\le x\, \}
=\{\, w\in I_{\phi}\mid w\le \phi(x)\, \}
\]
が成り立つ．
\end{cor}



\begin{lem}
$\phi$を$\rr$上の右連続な単調増加関数とする．
このとき任意の$y\in I_{\phi}$について
$y\le \phi(Q_{\phi}(y))$が成り立つ．
また，任意の$x\in \rr$について
$Q_{\phi}(\phi(x))\le x$が成り立つ．
\end{lem}
\begin{proof}
まず当然$Q_{\phi}(y)\le Q_{\phi}(y)$なので
先の補題から$y\le \phi(Q_{\phi}(y))$が成り立つ．
また，当然$\phi(x)\le \phi(x)$なので，当然
$Q_{\phi}(\phi(x))\le x$が成り立つ．
\end{proof}

\begin{lem}
$\phi$を$\rr$上の右連続な単調増加関数とする．
このとき
$Q_{\phi}$は右連続な単調増加関数
である．
\end{lem}
\begin{proof}
記述を簡単にするために$Q=Q_{\phi}$と書くことにする．
$Q$の単調性は定義から従う．

$w$を任意に固定する．
$\sup_{y<w}Q(w)=x$とおく．
すると
$x\le Q(w)$である．
$y<w$となる$y$を任意にとる．
今$x<Q(w)$と仮定して
$\epsilon=(Q(w)-x)/2$とする．
ここで，もしも$w\le \phi(x-\epsilon)$
とすると$Q(w)\le x-\epsilon<x$となるがこれは仮定に反するので$\phi(x-\epsilon)<w$である．
さらに，$Q(y)\le Q(w)$より，
$Q(y)+\epsilon\le x-\epsilon$が成り立つ．
よって
$y\le  \phi(Q(y))\le \phi(Q(y)+\epsilon)
\le 
\phi(x-\epsilon)<w$である．
ゆえに$w\le \phi(x-\epsilon)<w$となって矛盾する．
つまり$Q(w)=x$であり，$Q$の左連続性が示された．
\end{proof}

まず最初に$\phi$が有界な場合にLebesgue--Stieltjes測度を構成する.
\begin{thm}
$\phi:\rr\to \rr$を右連続な非減少関数とする．
さらに$\phi$は有界であると仮定する．
このとき，
$\rr$上の測度$\mu_{\phi}$で，
任意の$a, b\in [-\infty, \infty]$で$a<b$を満たすものに対して
\[
\mu_{\phi}((a, b])=\phi(b)-\phi(a)
\]
となるものが存在する．
\end{thm}
\begin{proof}
$I_{\phi}$上の通常のルベーグ測度を
$\nu$と書くことにする．
そして$\mu=(Q_{\phi})_{*}\nu$と書くことにしよう．
ここで$Q_{\phi}$は左連続だから特にボレル可測であることに注意しよう．
この$\mu$が求めている測度であることを以下で証明しよう．
まず任意の$x\in \rr$について
$\mu((-\infty, x])=
\nu(Q_{\phi}^{-1}((-\infty, x]))
$
であるが，
\[
Q_{\phi}^{-1}((-\infty, x])
=
\{\, w\in I_{\phi}\mid Q_{\phi}(w)\le x\, \}
=\{\, w\in I_{\phi}\mid w\le \phi(x)\, \}
\]
であるから，
$\mu((-\infty, x])=\phi(x)-\inf \phi$である．
$\phi$が有界なので，この値は有限の値であることに注意しよう．
よって
\[
\mu((a, b])=\phi(b)-\phi(a)
\]
となる．
\end{proof}

次に$\phi$が一般の場合にLebesgue--Stieltjes測度を構成しよう．
\begin{thm}
$\phi:\rr\to \rr$を右連続な非減少関数とする．
このとき，
$\rr$上の測度$\mu_{\phi}$で，
任意の$a, b\in [-\infty, \infty]$で$a<b$を満たすものに対して
\[
\mu_{\phi}((a, b])=\phi(b)-\phi(a)
\]
となるものが存在する．
\end{thm}
\begin{proof}
各$n$について$\phi_{n}i: \rr\to \rr$を
$\phi_{n}(s)=\min \{\max\{\phi(s), -n\}, n\}$と定義する．
これは
\[
\phi_{n}(s)=
\begin{cases}
n & \text{$n\le \phi(s)$のとき}\\
\phi(s) & \text{$-n\le \phi(s)<n$のとき}\\
-n & \text{$\phi(s)< -n$のとき}
\end{cases}
\]
と書いても同じことである．
この関数$\phi_{n}$は有界で右連続な単調関数になっている．
よって先の定理から
各$n$について
\[
\mu_{n}((a, b])=\phi_{n}(b)-\phi_{n}(a)
\]
となる測度$\mu_{n}$が存在する．
$\phi_{n}$は右連続なので二つの集合
$\{\, s\in \rr\mid n\le \phi(s)\, \}$
と
$\{\, s\in \rr\mid -n\le \phi(s)\, \}$
には最小元が存在する．
$\phi$は単調増加なので
$m_{n}\le M_{n}$である．
それをそれぞれ
$M_{n}$, $m_{n}$と書くことにする．
まず
任意の$n$と$A\in \mathfrak{B}(\rr)$について
$\mu_{n}(A)=\mu_{n}(A\cap [m_{n}, M_{n}])$が成り立つことを示そう．
この式が成り立つ$A$全体を$\mathfrak{S}$とおく．
すると$\mathfrak{S}$は$(a, b]$の形の集合を全て含んでいる．
$\mathfrak{S}$が$\lambda$系であることを示そう．
$X\in \mathfrak{S}$ $(\lambda1)$と
$\mathfrak{S}$が単調な族について閉じている$(\lambda2)$
ことは簡単にわかる．
$\mu_{n}(B\setminus A)=\mu_{n}(B)-\mu_{n}(A)=
\mu_{n}(B\cap [m_{n}, M_{n}])-\mu_{n}(A\cap [m_{n}, M_{n}])=
\mu_{n}((B\setminus A)\cap [m_{n}, M_{n}])$
なので$(\lambda 3)$も成り立つ．
よって$\pi$-$\lambda$定理によって，
$\mathfrak{S}=\mathfrak{B}(\rr)$となることがわかる．

同様にして，
任意の$n$と$A\in \mathfrak{B}(\rr)$について
\[
\mu_{n}(A\cap I_{n})=\mu_{n+1}(A\cap I_n)
\]
となることがわかる．

次に任意の$n$と$A\in \mathcal{B}(\rr)$について
$\mu_{n}(A)\le \mu_{n+1}(A)$
となることを示す．
まず$I_{n}=[m_{n}, M_{n}]$と置くと，定義から
$I_{n}\subset I_{n+1}$となることがわかる．
よって
$\mu_{n}(A)=\mu_{n}(A\cap I_n)=
\mu_{n+1}(A\cap I_n)
\le 
\mu_{n+1}(A\cap I_n)+\mu_{n+1}(A\cap (I_{n+1}\setminus I_{n}))
=
\mu_{n+1}(A\cap I_{n+1})$
なので成り立つ．

関数
$\mu: \mathfrak{B}(\rr)\to [0, \infty]$
を
$\mu(A)=\lim_{n\to \infty}\mu_{n}(A)$と定義する．
$\mu_{n}$の単調性からこの極限は存在する．
この$\mu$があの式を満たすことは$\phi_{n}$
が$\phi$に各点収束することからわかる．
$\mu$が測度であることを示そう．完全加法性を示せば良い．
$\sum$に関する単調収束定理，つまり数え上げ測度に関する
単調収束定理から
\begin{align*}
\mu\left(\coprod_{i\in \zz_{\ge 0}}A_{i}\right)
&=\lim_{n\to \infty}\mu_{n}\left(\coprod_{i\in \zz_{\ge 0}}A_{i}\right)
=
\lim_{n\to \infty}\sum_{i=0}^{\infty}\mu_{n}\left(A_{i}\right)
=
\sum_{i=0}^{\infty}\lim_{n\to \infty}\mu_{n}\left(A_{i}\right)\\
&=
\sum_{i=0}^{\infty}\mu\left(A_{i}\right)
\end{align*}
よって$\mu$は完全加法的である．
以上で一般の$\phi$に対しても
Lebesgue--Stieltjes測度が存在することがわかった．
\end{proof}


\section{おまけ：$\rr$についてのProkhorovの定理}
$\rr$に対するProkhovの定理はHellyの選出定理を使えば比較的簡単に示せるので，おまけとして紹介する．

以下の定理の証明は\cite{Helly}などを参考にしても良い．
\begin{thm}[Hellyの選出定理]\label{daiji3}
$\{F_n\}_{n\in \nn}$を$[0, 1]$に値をとる右連続な非減少関数の列とする．このとき部分列$\{F_{\phi(n)}\}$とある右連続な非減少関数$F$が存在して$F$の連続点$x\in \rr$毎に
\[
\lim_{n\to \infty}F_{\phi(n)}(x)=F(x)
\]
を満たす．
\end{thm}


\begin{thm}[$\rr$についてのProkhorovの定理]
$\mathcal{S}\subset \mathcal{P}(\rr)$とする．
このとき以下は同値である．
\begin{enumerate}
\item $\mathcal{S}$は相対コンパクトである．
\item $\mathcal{S}$は緊密である．
\end{enumerate}
\end{thm}
\begin{proof}
相対コンパクト性から緊密性を示す．
\[
	\ \forall \epsilon\in (0, 1)
	\ \exists m\in \nn
	\ \forall \mu\in \mathcal{S}\ 
	\mu((-m, m))\ge 1-\epsilon
\]
を示そう．
背理法による．
\[
	\ \exists \epsilon\in (0, 1)
	\ \forall m\in \nn
	\ \exists \mu\in \mathcal{S}\ 
	\mu((-m, m))< 1-\epsilon
\]
を仮定する．
各$m\in \nn$について対応する測度を$\mu_m$
とかく．
このとき，$\mathcal{S}$のコンパクト性から，
部分列$\{m_i\}$と$\{\mu_i\}$
が存在して$x\le m_i$ならば
\[
\mu_i((-x, x))<1-\epsilon
\]
と
\[
\mu_i\To \mu
\]
が成り立つ．
そして弱収束していることと$m_i\to \infty$から
任意の$x\in \rr$について
\[
\mu((-x, x))\le \liminf_{i\to \infty}\mu_i((-x. x))\le 1-\epsilon
\]
が成り立つが，$x\to \infty$とすると，
$1\le 1-\epsilon$となるので矛盾する．

今度は逆を示そう．
$\mathcal{S}$
は緊密であるとする．このとき
\[
	\forall N\in \nn\ \exists M_N\in \nn\ \forall \mu\in \mathcal{S}\ 
	\mu([-M_N, M_N])\ge 1-2^{-N}\land M_N\ge N
\]
が成り立つ．
$\mathcal{S}$内の部分列$\{\mu_n\}_{n\in \nn}$をとる．
$\mu_n$の分布関数を$F_n$と書くことにする．
このとき，
Hellyの選出定理から
$F_n$の部分列と単調非減少関数$F$で，
$F$の連続点について
各点収束する．
緊密性から，$M_N\le x$を満たす$x$について
\[
1\ge F_n(x)-F_n(-x)\ge 1-2^{-N}
\]
であり，$F$の不連続点は高々可算個しかないので，
$x$を連続点として$n\to \infty$とすると，
\[
1\ge F(x)-F(-x)\ge 1-2^{-N}
\]
が成り立つ．$x\to \infty$
とすると，
\[
\lim_{x\to \infty}(F(x)-F(-x))=1
\]
が成り立つ．
また，$\lim_{x\to \infty}F(-x)=0$は$F_n$の収束先であることからすぐわかるので，
この文章の前半で頑張って示した内容から
$F$はある確率測度$\mu$の分布関数になる．
つまり，$\mu_n\To \mu$となるので，$\mathcal{S}$が相対コンパクトであることがわかる．
\end{proof}










	
\begin{thebibliography}{9}
\bibitem{Helly}
はてなブログ電波通信の記事
「Helly空間とHellyの選出定理」,
https://concious4410.hatenablog.com/entry/2021/03/25/100947


\bibitem{RMK1}
はてなブログ電波通信の記事
「リース・マルコフ・角谷の表現定理の証明(入射角16.225°の測度論入門2)」, 
https://concious4410.hatenablog.com/entry/2015/10/28/213025

\bibitem{RMK2}
はてなブログ電波通信の記事
「Riesz-Markov-Kakutaniの表現定理II」, 
https://concious4410.hatenablog.com/entry/2020/02/02/213341


\bibitem{Ito}伊藤清三，「ルベーグ積分入門」，
1963年，
裳華房．

\bibitem{Hu}船木直久, 「確率論」, 
講座・数学の考え方 20，2004, 朝倉書店．
\end{thebibliography}




			\end{document}



\begin{comment}
[集合のやつは$\mid$を使う。]
[真なる包含関係]
\subsetneqq
[可換図式]
\[\xymatrix{
&   & (2) \ar@{=>}[rd]&  &  &  & (5) \ar@{=>}[rd]&\\ 
& (1)\ar@{=>}[ru] \ar@{=>}[rd]&  &(4)\ar@{=>}[r]&(7)& (1)\ar@{=>}[ru] \ar@{=>}[rd]& & (7)&\\ 
&   & (3) \ar@{=>}[ur]&  &  &  &(6) \ar@{=>}[ur]&
}\]

[\bigin \endを使わない方法]

{\prop[aa] aaa}
で
\begin{prop}[aa]
aaa
\end{prop}
と同じ\df \thm \proof でも同様

[直和]
$\amalg$

[同値関係でよく使う波線]
\sim

[引数付きの新命令]
\newcommand{\prn}[1]{\|#1\|_p}

[番号のリセット]

\setcounter{equation}{0}


[字体の変更]

{\rm hogehoge}と使う。

\rm ローマン
\it イタリック
\sl 斜体

[supp]

\newcommand{\supp}{\text{{\rm supp}}}

[中央揃えの改行]

	\begin{eqnarray*}
		hoge1&=&hogehoge
		hoge2&=&hogehofe

	\end{eqnarray*}

&&の部分で揃えられる。


[\tag]
\begin{equation}
f(x) = ax^2 + bx + c \tag{1.2.3}
\end{equation}



[場合分け]

	\begin{cases}
		hoge1 & \text{状況１}\\
		hoge2 & \text{状況２}\\
		・
		・
		・
	\end{cases}



\end{comment}





